\section{Proofs of the composite and prediction claims}
\label{hfpr:appendix}

This appendix proves the mathematical implications used in this report.
The hypotheses are part of each result: a proof under a Gaussian model
does not prove that model for human scores. The foundations are linearity
of expectation and conditional expectation, covariance expansion and the
Gaussian characteristic function \citep[Sections 1.6, 3.3 and 4.1]{durrett2019},
Gaussian affine closure and regression \citep[Theorems 3.3.1--3.3.2 and
3.5.1--3.5.2]{gallager2013}, and the spectral theorem for real symmetric
matrices \citep[Theorem 7.29]{axler2024}. Their relevant forms and proofs
are stated in the shared mathematical foundations preceding this appendix.
In particular, Propositions~\ref{pf:moments}, \ref{pf:spectral},
\ref{pf:gaussian}, \ref{pf:conditioning} and \ref{pf:expectation}
state the moment, spectral, Gaussian and conditional-expectation tools.
Write \(\phi(x)=(2\pi)^{-1/2}\exp(-x^2/2)\),
\(\Phi(x)=\int_{-\infty}^x\phi(t)\,dt\), and
\(\overline\Phi=1-\Phi\).

\begin{proposition}[Covariance, scales and weighted composites]
\label{hfpr:standardization}
For square-integrable scores \(X_i\) with means \(\mu_i\) and positive
SDs \(\sigma_i\), put \(Z_i=(X_i-\mu_i)/\sigma_i\). Then
\(\E Z_i=0\), \(\Var(Z_i)=1\), and their covariance matrix is their
correlation matrix \(R\). For any fixed real vector \(w\),
\[
 \E(w^{\mathsf T}Z)=0,\qquad
 \Var(w^{\mathsf T}Z)=w^{\mathsf T}Rw.
\]
When the latter is positive, division by its square root gives a mean-zero,
unit-variance composite. The transformation \(I=100+sZ_C\), \(s>0\),
has population mean 100 and SD \(s\).
\end{proposition}
\begin{proof}
Linearity gives \(\E(X_i-\mu_i)=0\). Expanding the square in the
definition of variance gives \(\Var(aX+b)=a^2\Var(X)\); this proves
the first two conclusions and, on expanding a product, gives
\(\Cov(Z_i,Z_j)=\Cov(X_i,X_j)/(\sigma_i\sigma_j)=r_{ij}\).
For any finite collection, integrability of each product follows from
Cauchy--Schwarz. Hence finite sums can be taken through expectation:
\[
 \Var\!\left(\sum_iw_iZ_i\right)
 =\E\!\left[\sum_{i,j}w_iw_jZ_iZ_j\right]
 =\sum_{i,j}w_iw_jr_{ij}=w^{\mathsf T}Rw.
\]
For two unstandardized scores the same centered expansion gives
\(\Var(X+Y)=\Var(X)+\Var(Y)+2\Cov(X,Y)\).
The assertions about division and reporting scales now follow by the
same affine-variance identity. No distributional assumption was used.
For example, \((160-100)/15=4\) and \((160-100)/16=15/4=3.75\).
\end{proof}

\begin{proposition}[Equal weights, admissibility and information in the matrix]
\label{hfpr:equal-weights}
For \(n\geq2\), let \(\bar r=2\sum_{i<j}r_{ij}/[n(n-1)]\). Then
\[
 v_n:=\Var(\bar Z)=\frac{1+(n-1)\bar r}{n}
 =\bar r+\frac{1-\bar r}{n},\qquad
 -\frac1{n-1}\leq\bar r\leq1.
\]
When \(v_n>0\), the composite is \(\bar Z/\sqrt{v_n}\).
The entire correlation matrix must be positive semidefinite (PSD);
the displayed bound on its average is necessary, rather than a sufficient
test for an arbitrary proposed matrix. For an equicorrelation matrix the
same bound is necessary and sufficient for PSD. The average correlation
alone need not determine the variance with unequal weights.
\end{proposition}
\begin{proof}
Apply Proposition~\ref{hfpr:standardization} with every weight \(1/n\).
The diagonal contributes \(n/n^2\); symmetry makes the off-diagonal
contribution \(2\sum_{i<j}r_{ij}/n^2\). Substitution proves the formula.
Every quadratic form \(u^{\mathsf T}Ru=\Var(u^{\mathsf T}Z)\) is
nonnegative, proving PSD and, with \(u\) the all-ones vector, the lower
bound on \(\bar r\). Cauchy--Schwarz gives \(|r_{ij}|\leq1\), proving
the upper bound. For \(n=1\), directly \(v_1=1\).

In the equicorrelated case \(R=(1-c)I+c\mathbf1\mathbf1^{\mathsf T}\).
Its action on \(\mathbf1\) multiplies by \(1+(n-1)c\); its action
on every vector orthogonal to \(\mathbf1\) multiplies by \(1-c\).
The spectral theorem therefore gives PSD exactly for
\(-1/(n-1)\leq c\leq1\). Such a PSD matrix is a realizable Gaussian
covariance: diagonalize it, take nonnegative square roots of its
eigenvalues, and apply the resulting square-root matrix to independent
standard normals. A zero \(v_n\) makes the mean deterministic and the
division undefined.

To prove the unequal-weight qualification, compare \(R_0=I_3\) with
\[
 R_t=\begin{pmatrix}1&t&-t\\t&1&0\\-t&0&1\end{pmatrix},\qquad t=1/4.
\]
Both average off-diagonal correlations are zero. The latter matrix has
eigenvalues \(1,1+\sqrt2t,1-\sqrt2t\), all positive, as is checked
by its characteristic polynomial
\((1-\lambda)[(1-\lambda)^2-2t^2]\).
For \(w=(1,1,0)^{\mathsf T}\), however,
\(w^{\mathsf T}R_tw=2+2t\ne2=w^{\mathsf T}R_0w\).
Finally, if \(I_i=100+sZ_i\), then
\(\bar Z=(\bar I-100)/s\). Substitution into
\(100+s\bar Z/\sqrt{v_n}\) proves the IQ-style formula.
\end{proof}

\begin{proposition}[Gaussian percentiles require a joint model]
\label{hfpr:normality}
If \(Z\) is jointly Gaussian and \(w^{\mathsf T}Rw>0\), its
standardized fixed weighted composite is standard normal and its
model-based percentile is \(100\Phi(z_C)\). Neither two moments nor
normal marginal distributions imply this conclusion.
\end{proposition}
\begin{proof}
An affine image of a jointly Gaussian vector is Gaussian
\citep[Theorem 3.3.1]{gallager2013}; Proposition~\ref{hfpr:standardization}
gives its mean zero and variance one. By definition of its CDF, the fraction
at or below \(z_C\) is then \(\Phi(z_C)\).

For the failure of the converse, take \(X\sim N(0,1)\) and an independent
sign \(S\) with \(P(S=1)=(1+c)/2\) and \(P(S=-1)=(1-c)/2\),
where \(-1<c<1\). Set \(Y=SX\). Symmetry of the normal distribution
makes \(Y\) standard normal. Independence gives
\(\Cov(X,Y)=\E(SX^2)=\E S\,\E X^2=c\).
Nevertheless \(P(X+Y=0)=(1-c)/2>0\): on \(S=-1\) the sum vanishes,
whereas on \(S=1\) it is \(2X\), which has no atom. A nondegenerate
normal distribution has a density and hence no atoms. Standardization
preserves this atom, so these normal marginals and their correlation do
not imply a normal composite. They also fail to identify joint tail
probabilities: for positive \(a,b\), this model gives
\(P(X>a,Y>b)=(1+c)\overline\Phi(\max(a,b))/2\), since opposite signs
cannot both exceed positive cutoffs. It is not the Gaussian joint law.
The same counterexample, having the required moments, proves that moments
alone do not determine the tail.
\end{proof}

\begin{proposition}[Thresholds, monotonicity and fixed-correlation limits]
\label{hfpr:thresholds}
Let \(q=\Phi^{-1}(1-10^{-6})\) be defined by the continuous standard
normal law. For a Gaussian composite with mean-variance factor \(v_n>0\),
the one-in-a-million threshold on the component mean is
\(t_n(c)=q\sqrt{c+(1-c)/n}\). For \(0\leq c\leq1\) it decreases
with \(n\), increases with \(c\) for \(n>1\), and tends to
\(q\sqrt c\). For fixed positive observed mean, its standardized
composite decreases with \(c\). These are properties of the stipulated
sequence of batteries, rather than a bound on possible score values.
\end{proposition}
\begin{proof}
Since \(\phi>0\), integration shows that \(\Phi\) is continuous and
strictly increasing from 0 to 1. The quantile is uniquely defined and
\(P(Z_C>q)=1-\Phi(q)=10^{-6}\). Because \(\sqrt{v_n}>0\),
\(\bar Z>q\sqrt{v_n}\) is exactly \(Z_C>q\), proving the cutoff.
The normal law has no atoms, so strict or non-strict inequalities have
the same model probability.

The ordering \(q>3>2\), used for nested simulated counts, can be shown
without relying on its printed decimal. On \([3,4]\) the density is
at least \(\phi(4)\), and
\[
 \overline\Phi(3)\geq\int_3^4\phi(x)\,dx\geq\phi(4)
 =\frac{1}{\sqrt{2\pi}e^8}>
 \frac{1}{3\cdot3^8}=\frac1{19683}>10^{-6}.
\]
Here \(e=\sum_{j=0}^{\infty}1/j!<3\), by comparison of terms
\(j!\geq2^{j-1}\) for \(j\geq1\) with strict inequality for
\(j\geq3\), and \(\pi=4\int_0^1(1+x^2)^{-1}\,dx<4\), so
\(\sqrt{2\pi}<3\). Strict decrease of the normal upper tail proves
\(q>3\).

For real \(x\geq1\) put \(v(x,c)=c+(1-c)/x\). Elementary
differentiation gives
\[
 \frac{\partial t}{\partial x}
 =-\frac{q(1-c)}{2x^2\sqrt{v(x,c)}},\qquad
 \frac{\partial t}{\partial c}
 =\frac{q(1-1/x)}{2\sqrt{v(x,c)}}.
\]
The signs prove the claims, with equality in the evident singleton or
perfect-correlation cases. For a fixed \(b>0\),
\(\partial[b/\sqrt{v(x,c)}]/\partial c
=-b(1-1/x)/(2v(x,c)^{3/2})\leq0\).
For \(c<1\), the raw-SD derivative has magnitude
\((1-c)/(2\sqrt{x^3(cx+1-c)})\). The positive denominator increases
with \(x\geq1\) for \(0\leq c<1\), proving the slowing decline;
for \(c=1\) it is constant.
Also \(v_n-c=(1-c)/n\to0\); continuity of square root gives the limit.
When \(c>0\), rationalizing yields the explicit error bound
\[
 0\leq t_n(c)-q\sqrt c
 =\frac{q(1-c)}{n(\sqrt{v_n}+\sqrt c)}
 \leq\frac{q(1-c)}{2n\sqrt c}.
\]
For \(c=0\) the exact cutoff is \(q/\sqrt n\to0\). Fixed negative
\(c\) cannot give valid equicorrelated batteries for arbitrarily large
\(n\), by Proposition~\ref{hfpr:equal-weights}. Under the Gaussian model
every nondegenerate \(Z_C\) has unbounded support; thus a limit on the
component-mean \emph{cutoff} is not a ceiling on attainable composites.
Changing the correlations as tests are added changes the sequence and
does not satisfy the fixed-\(c\) hypothesis. Finally \(100+s q\) is
the composite reporting-scale cutoff for every \(n\), whereas
\(100+s t_n(c)\) is the component-average cutoff, by the affine identity.
\end{proof}

\begin{proposition}[Common-factor aggregate, posterior and residual contrasts]
\label{hfpr:factor}
Assume \(0<c<1\) and
\(X_i=\sqrt c\,G+\sqrt{1-c}\,\varepsilon_i\), where all displayed
basic variables are independent standard normals. For
\(M_n=n^{-1}\sum_iX_i\), put \(v_n=c+(1-c)/n\) and
\(R_n=c/v_n=nc/[1+(n-1)c]\). Then
\[
 (G,M_n)\ \stackrel d=\
 \left(G,\sqrt c\,G+\sqrt{(1-c)/n}\,E\right),
\]
where \(E\) is standard normal independent of \(G\). Given the entire
score vector, \(G\) is normal with mean
\(\widehat G_n=(\sqrt c/v_n)M_n=\sqrt{R_n}\,Z_{C,n}\) and variance
\(1-R_n=(1-c)/[1+(n-1)c]\). Moreover
\[
 \operatorname{Corr}(G,Z_{C,n})^2=R_n,\quad
 \Var(\widehat G_n)=R_n,\quad
 \E[(G-\widehat G_n)^2]=1-R_n.
\]
The conditional mean minimizes mean-square prediction error among all
square-integrable functions of the scores.
\end{proposition}
\begin{proof}
Linearity, independence and unit variances give \(\E X_i=0\),
\(\Var(X_i)=1\), \(\Cov(X_i,X_j)=c\) for \(i\ne j\), and
\(\Cov(G,M_n)=\sqrt c\). Define
\(E=n^{-1/2}\sum_i\varepsilon_i\). Its Gaussian affine law has mean
zero and variance one, and it is independent of \(G\) because it is a
function of residuals independent of \(G\). Substitution gives the
stated representation, including its joint law with \(G\).

Let \(W=G-(\sqrt c/v_n)M_n\). Its variance, obtained by expanding
the square, is \(1+c/v_n-2c/v_n=1-c/v_n\), and
\(\Cov(W,M_n)=\sqrt c-(\sqrt c/v_n)v_n=0\). For each contrast
\(D_i=X_i-M_n\), symmetry gives
\[
 \Cov(G,D_i)=\sqrt c-\sqrt c=0,\qquad
 \Cov(M_n,D_i)=\frac{1+(n-1)c}{n}-v_n=0.
\]
Thus \(\Cov(W,D_i)=0\) as well. All these variables are jointly
Gaussian. Their characteristic function is
\(\exp(it^{\mathsf T}\mu-t^{\mathsf T}\Sigma t/2)\); zero
cross-covariance makes it factor between \(W\) and \((M_n,D_1,\ldots,D_n)\).
Uniqueness of characteristic functions proves independence even though
the contrasts are linearly dependent. This is the Gaussian regression
principle \citep[Theorems 3.3.2 and 3.5.1--3.5.2]{gallager2013}.
Since the score vector is recovered by \(X_i=M_n+D_i\), \(W\) is
independent of all observed scores. Conditional on them,
\(G=(\sqrt c/v_n)M_n+W\), with the stated normal mean and variance.
This explicitly proves the claim that the contrasts add no information
about \(G\) under this symmetric Gaussian model.

The covariance with \(Z_{C,n}=M_n/\sqrt{v_n}\) is
\(\sqrt{c/v_n}\). Both \(G\) and \(Z_{C,n}\) have variance one,
proving the squared-correlation formula. Scaling \(M_n\) gives
\(\Var(\widehat G_n)=c/v_n\). The centered independent residual gives
the MSE formula. For any square-integrable score function \(h\),
conditional expectation and pull-out give
\(\E[W(\widehat G_n-h)]=0\), since \(\E[W\mid X]=0\).
Expanding \(G-h=W+(\widehat G_n-h)\) proves
\[
 \E[(G-h)^2]=1-R_n+\E[(\widehat G_n-h)^2]\geq1-R_n.
\]
Finally \(1-R_n=(1-c)/(nc+1-c)\to0\) for fixed \(c>0\), although
\(\SD(M_n)=\sqrt{v_n}\to\sqrt c>0\). Population spread of scores
and conditional prediction error are thus different mathematical objects.
\end{proof}

\begin{proposition}[Observed score laws do not identify a latent interpretation]
\label{hfpr:nonidentification}
Even an exactly known observable correlation matrix, or the entire joint
score law, need not identify a proposed latent ability, its posterior, or
an individual measurement-error interval.
\end{proposition}
\begin{proof}
Keep precisely the same observable \(X\) from
Proposition~\ref{hfpr:factor}. In one model call its common factor \(G\)
the target. In a second model extend the probability space by a standard
normal \(H\) independent of \((G,X)\), and call \(H\) the target.
Every observable distribution and its matrix \(R\) are unchanged, yet
\(\E[H\mid X]=0\), \(\Var(H\mid X)=1\), whereas the first model has
the different posterior proved in Proposition~\ref{hfpr:factor}.
Therefore no rule depending only on those observable quantities can
deduce both latent posteriors correctly. Additional measurement assumptions
or external validity evidence are logically required. The example does not
claim that either target is the correct psychometric construct.
\end{proof}

\begin{proposition}[Selection and repeated-score procedures can alter calibration]
\label{hfpr:selection}
Population correlations need not equal selected-sample correlations.
Likewise the maximum or a stopping rule based on repeated attempts need
not have the distribution of one prespecified score. The direction and
size of an adjustment require assumptions about selection and dependence.
\end{proposition}
\begin{proof}
Let independent \(X,E\) be standard normal and let
\(Y=\rho X+\sqrt{1-\rho^2}\,E\), \(0<\rho<1\).
Unconditionally \(\operatorname{Corr}(X,Y)=\rho\).
Select on \(X>0\). Symmetry gives \(\E[X^2\mid X>0]=1\), while
\[
 \E[X\mid X>0]=2\int_0^\infty x\phi(x)\,dx
 =2\phi(0)=\sqrt{2/\pi},
\]
using \(\phi'(x)=-x\phi(x)\). Consequently
\(v=\Var(X\mid X>0)=1-2/\pi\in(0,1)\).
Independence of \(E\) from the selection event gives conditional
covariance \(\rho v\) and variance of \(Y\) equal to
\(\rho^2v+1-\rho^2\). The selected correlation is therefore
\[
 \rho\sqrt{\frac{v}{\rho^2v+1-\rho^2}}<\rho,
\]
because \(\rho^2v+1-\rho^2-v=(1-\rho^2)(1-v)>0\).
No selection gives the original value, so there is already no universal
correction independent of the mechanism. This constructive mathematical
example supports the possibility of distortion; empirical correction
methods are attributed to \citet{sackett2000}.

For independent standard-normal attempts \(U_1,\ldots,U_m\),
independence gives
\(P(\max_jU_j\leq t)=\Phi(t)^m\), so the chance of any qualification
above a fixed cutoff is \(1-\Phi(t)^m\), exceeding the one-attempt chance
for finite \(t\) and \(m>1\). The probability of no success after \(m\)
attempts tends to zero, proving eventual qualification with probability
one if attempts can continue independently without limit. In contrast,
perfectly repeated scores \(U_j=U_1\) give qualification probability
\(1-\Phi(t)\) for every \(m\). These are explicit counterexamples to
treating all procedures as interchangeable; no claim that every possible
dependent procedure differs is needed. Practice effects themselves are
an empirical finding attributed to \citet{hausknecht2007}.

Finally a finite number \(m\) of observations cannot logically determine
an unknown rare-event probability: the event of no exceedances has
probability \((1-p)^m>0\) for every \(0<p<1\) in an independent
Bernoulli model. At \(m=27\) and \(p=10^{-6}\), the union bound gives
\(P(\hbox{no event})\geq1-27\cdot10^{-6}=0.999973\).
Selection further changes the sampling experiment. A small selected
sample is therefore not a deductive establishment of a one-in-a-million
population frequency.
\end{proof}

\begin{proposition}[Coupling, rare-event precision and an exponential bound]
\label{hfpr:simulation}
The aggregate simulation preserves each fixed-size \((G,M_n)\) law but
reusing a residual across sizes does not simulate a nested battery.
Nested exceedance indicators imply nested counts. Under independent
ideal draws with event probability \(p>0\), relative Monte Carlo standard
error is \(\sqrt{(1-p)/(Tp)}\). For \(0<p<1\),
\[
 e^{-Tp-Tp^2/[2(1-p)]}\leq(1-p)^T\leq e^{-Tp}.
\]
\end{proposition}
\begin{proof}
The aggregate representation is proved in Proposition~\ref{hfpr:factor}.
With one shared \(E\), two sizes \(m<n\) have covariance
\(c+(1-c)/\sqrt{mn}\). In a genuine nested battery their residual
averages share only the first \(m\) terms, giving instead
\(c+(1-c)/n\), by covariance expansion. These differ for \(0<c<1\),
so the fixed-size shortcut is not a claim about nested paths. As a concrete
dependence diagnostic, two aggregates formed with the shared \((G,E)\)
have positive covariance; they cannot be independent centered variables.
The per-size laws still follow from the separately proved representation.

For \(q>3>2\), every real \(z\) satisfies
\(\mathbf1_{z>q}\leq\mathbf1_{z>3}\leq\mathbf1_{z>2}\).
Summing pointwise proves the event-count ordering without any sampling
assumption. For independent Bernoulli indicators, covariance expansion
gives \(\Var(K/T)=p(1-p)/T\); divide its SD by \(p\) to obtain the
relative error. For \(p=10^{-6}\) and \(T=4\cdot10^6\) it is
\(\sqrt{(1-10^{-6})/4}\), approximately one half. A target relative
SE at most \(\delta\) requires exactly
\(T\geq(1-p)/(p\delta^2)\); at \(\delta=0.1\) this is
\(T\geq99{,}999{,}900\). This is a standard-error calculation, not a
95\% confidence-width guarantee.

Independence gives \(P(K=0)=(1-p)^T\). The fundamental theorem of
calculus yields
\[
 -\log(1-p)-p=\int_0^p\frac{t}{1-t}\,dt,
 \qquad0\leq\int_0^p\frac{t}{1-t}\,dt
 \leq\frac{p^2}{2(1-p)}.
\]
Multiply by \(T\), exponentiate and preserve the inequalities to obtain
the stated bound. For \(p=10^{-6}\) and \(T=10^6\) or \(4\cdot10^6\),
the exponent correction is at most \(5.000006\cdot10^{-7}\) or
\(2.000003\cdot10^{-6}\), respectively; hence the quoted
\(e^{-1}\) and \(e^{-4}\) approximations have explicitly controlled
relative error. Binomial count moments, Gaussian moment diagnostics,
pooled variance and coverage of the reported intervals are proved in the
shared statistical-reference appendix (Propositions~\ref{sm:normal-sample},
\ref{sm:binomial} and \ref{sm:pooling}, and the interval results
\ref{sm:sd-interval}, \ref{sm:wilson}, \ref{sm:clopper-pearson}); those laws assume ideal independent
draws, not merely a seeded finite-state generator.
For the empirical factor MSE specifically, the independent residuals across
draws obey \(W_j/\sqrt{1-R_n}\sim N(0,1)\), by
Proposition~\ref{hfpr:factor}. Therefore
\(\sum_{j=1}^T W_j^2/(1-R_n)\) is chi-square with \(T\) degrees of
freedom by its defining sum-of-squares law. This also identifies its
expectation \(T\) and thus the empirical MSE target \(1-R_n\).
\end{proof}

\begin{proposition}[Exact arithmetic behind the illustrative numbers]
\label{hfpr:numerics}
For three SD-16 scores averaging 160 and \(c=3/5\),
\[
 v_3=\frac{11}{15},\qquad
 Z_C=\sqrt{\frac{3375}{176}},\qquad
 I_C=100+\sqrt{\frac{54000}{11}}.
\]
For the latent example when \(Z_C=q\) at the exact quantile,
\[
 R_3=\frac9{11},\qquad
 \E[G\mid X]=\frac3{\sqrt{11}}q,\qquad
 \SD(G\mid X)=\sqrt{\frac2{11}}.
\]
The threshold table consists exactly of \(q\sqrt{(3n+2)/(5n)}\)
and its affine SD-16 transform; its infinite row is \(q\sqrt{3/5}\).
\end{proposition}
\begin{proof}
Substitute \(\bar z=15/4\), \(c=3/5\), \(n=3\) into
Propositions~\ref{hfpr:equal-weights} and \ref{hfpr:factor}. Squaring
the positive worked composite gives
\((15/4)^2/(11/15)=3375/176\); squaring its SD-16 excess gives
\(16^2\cdot3375/176=54000/11\). Exact rational comparisons give
\[
 (4.37905)^2<3375/176<(4.37915)^2,\qquad
 (70.055)^2<54000/11<(70.065)^2.
\]
Monotonicity of the positive square root proves rounding to \(4.3791\)
and \(170.06\). For the posterior,
\(R_3=(9/5)/(11/5)=9/11\) and \(1-R_3=2/11\), proving the formulas.
The rational inequalities
\((0.425)^2<2/11<(0.435)^2\) and
\((6.815)^2<512/11<(6.825)^2\) prove the displayed conditional-SD
roundings \(0.43\) and \(6.82\).
The scale conversion is \(100+16(3/\sqrt{11})q\) for its mean and
\(16\sqrt{2/11}\) for its SD. The SD-15 singleton threshold is
\(100+15q\); the three-test threshold at \(c=9/10\) is
\(100+16q\sqrt{14/15}\). These exact expressions account for every
worked decimal and threshold-table entry.

The displayed quantile and its resulting decimal values are numerical
approximations, reproduced by the documented R upper-tail quantile or
\texttt{scipy.stats.norm.isf} and the repository's calculation script.
Agreement of these software evaluations is numerical evidence, not by
itself an exact-real proof of their final rounded digits. Sensitivity to
quantile error is explicit: replacing \(q\) by \(q+\eta\) changes a
component threshold by exactly \(\eta\sqrt{v_n}\), its SD-\(s\) version
by \(s\eta\sqrt{v_n}\), and the latent-scale mean by
\(16\sqrt{R_n}\eta\). Thus the algebra is proved independently of
backend accuracy, and any certified quantile enclosure can be propagated
by these positive factors. Likewise recorded sample means, variances,
counts and interval endpoints are observations/computed approximations,
not identities asserted of all realizations.
\end{proof}
